\documentclass{../../nndlcourse}

\begin{document}

\title{实验4 - 循环神经网络}
\author{数据科学与大数据技术专业}
\date{第9周}
\maketitle


% ============================================================
\section{实验目标}
% ============================================================
\begin{itemize}
    \item 理解循环神经网络(RNN)如何通过隐藏状态处理序列数据.
    \item 使用 NumPy 实现基础 RNN 单元、完整 RNN 前向传播、LSTM 单元和完整 LSTM 前向传播.
    \item 理解通过时间反向传播(BPTT)、梯度爆炸、梯度裁剪和序列生成的基本机制.
    \item 完成恐龙名字字符级语言模型, 能够训练 RNN 并采样生成新名字.
    \item 阅读并复现 Karpathy 的 \texttt{minimalRNN.py}, 理解字符级 Vanilla RNN 的完整训练闭环.
    \item 掌握 one-hot 输入、softmax 交叉熵、隐藏状态续接、AdaGrad 参数更新和采样生成之间的关系.
\end{itemize}


% ============================================================
\section{实验环境}
% ============================================================
\begin{itemize}
    \item 操作系统: Windows / Linux / macOS
    \item Python 3.7 及以上
    \item 主要依赖库: \texttt{numpy}、\texttt{matplotlib}
    \item 推荐开发环境: Jupyter Notebook 或 VS Code
    \item 实验文件:
    \begin{itemize}
        \item \texttt{W1A1/Building\_a\_Recurrent\_Neural\_Network\_Step\_by\_Step.ipynb}
        \item \texttt{W1A2/Dinosaurus\_Island\_Character\_level\_language\_model.ipynb}
        \item \texttt{2025-实验4.ipynb}
        \item \texttt{minimalRNN.py}
    \end{itemize}
\end{itemize}

如本地环境缺少依赖, 可在终端中执行:
\begin{verbatim}
pip install numpy matplotlib
\end{verbatim}


% ============================================================
\section{背景知识}
% ============================================================

\subsection{序列建模与隐藏状态}

循环神经网络用于处理文本、时间序列、语音等有顺序的数据. 与前馈网络不同, RNN 在每个时间步共享同一组参数, 并通过隐藏状态保存历史信息. 对时间步 $t$, 基础 RNN 的递推为:
\begin{align}
    a^{\langle t\rangle}
    &= \tanh\left(W_{aa}a^{\langle t-1\rangle}
    + W_{ax}x^{\langle t\rangle} + b_a\right) \\
    \hat{y}^{\langle t\rangle}
    &= \operatorname{softmax}\left(W_{ya}a^{\langle t\rangle} + b_y\right)
\end{align}
其中 $x^{\langle t\rangle}$ 是当前输入, $a^{\langle t\rangle}$ 是隐藏状态, $\hat{y}^{\langle t\rangle}$ 是当前预测.

本实验约定输入序列张量形状为:
\begin{align}
    x \in \mathbb{R}^{n_x\times m\times T_x}
\end{align}
其中 $n_x$ 为输入维度, $m$ 为 batch 大小, $T_x$ 为序列长度.

\subsection{LSTM 门控机制}

基础 RNN 对长程依赖建模能力有限, 训练时容易出现梯度消失或梯度爆炸. LSTM 通过细胞状态和门控机制缓解这一问题. 单步 LSTM 的主要公式为:
\begin{align}
    \Gamma_f &= \sigma(W_f[a^{\langle t-1\rangle};x^{\langle t\rangle}] + b_f) \\
    \Gamma_i &= \sigma(W_i[a^{\langle t-1\rangle};x^{\langle t\rangle}] + b_i) \\
    \tilde{c}^{\langle t\rangle} &= \tanh(W_c[a^{\langle t-1\rangle};x^{\langle t\rangle}] + b_c) \\
    c^{\langle t\rangle} &= \Gamma_f\odot c^{\langle t-1\rangle}
    + \Gamma_i\odot \tilde{c}^{\langle t\rangle} \\
    \Gamma_o &= \sigma(W_o[a^{\langle t-1\rangle};x^{\langle t\rangle}] + b_o) \\
    a^{\langle t\rangle} &= \Gamma_o\odot \tanh(c^{\langle t\rangle})
\end{align}
遗忘门决定保留多少旧记忆, 输入门决定写入多少新信息, 输出门决定暴露多少细胞状态给隐藏状态.

\subsection{字符级语言模型}

字符级语言模型以单个字符为基本 token. 给定字符序列 $x_1,\ldots,x_T$, 模型学习:
\begin{align}
    P(x_1,\ldots,x_T)=\prod_{t=1}^{T}P(x_t\mid x_1,\ldots,x_{t-1})
\end{align}
训练时, 输入为当前字符, 目标为下一个字符. 每个字符通常转为 one-hot 列向量, 输出层使用 softmax 得到下一个字符的概率分布.

对一段序列, 交叉熵损失为:
\begin{align}
    \mathcal{L} = -\sum_{t=1}^{T}\log \hat{p}^{\langle t\rangle}_{y^{\langle t\rangle}}
\end{align}
推理时, 模型从起始输入开始, 每一步根据 softmax 概率采样一个字符, 再把该字符作为下一步输入, 循环生成文本.

\subsection{BPTT 与梯度裁剪}

RNN 的反向传播需要沿时间轴从后向前展开, 简称 BPTT. 因为隐藏状态递归依赖过去, 当前时间步的损失会影响前面许多时间步的参数梯度.

长序列训练时梯度可能指数级增大. 梯度裁剪将梯度限制在固定范围内:
\begin{align}
    g \leftarrow \operatorname{clip}(g, -\tau, \tau)
\end{align}
本实验中恐龙名字模型和 minimalRNN 补充实验都使用裁剪阈值 5.

\subsection{AdaGrad 参数更新}

\texttt{minimalRNN.py} 使用 AdaGrad 更新参数. 对参数 $\theta$ 和梯度 $g$:
\begin{align}
    G &\leftarrow G + g^2 \\
    \theta &\leftarrow \theta - \eta \frac{g}{\sqrt{G+\epsilon}}
\end{align}
其中 $G$ 是历史平方梯度累积缓存. AdaGrad 会自动缩小频繁更新参数的有效学习率, 有助于稳定小型 RNN 训练.


% ============================================================
\section{实验步骤}
% ============================================================

\subsection*{实验4.1: RNN 与 LSTM 前向传播(W1A1)}

\textbf{实验目标: }使用 NumPy 手写基础 RNN 和 LSTM 的前向传播, 明确每个时间步的输入、隐藏状态、输出和 cache.

\subsubsection*{4.1.1 RNN 单元}

\textbf{练习 1: \texttt{rnn\_cell\_forward}}

实现单步 RNN 前向传播:
\begin{align}
    a_{\mathrm{next}} &= \tanh(W_{aa}a_{\mathrm{prev}} + W_{ax}x_t + b_a) \\
    y_t &= \operatorname{softmax}(W_{ya}a_{\mathrm{next}} + b_y)
\end{align}
函数返回 \texttt{a\_next}、\texttt{yt\_pred} 和 \texttt{cache=(a\_next,a\_prev,xt,parameters)}.

\subsubsection*{4.1.2 完整 RNN 前向传播}

\textbf{练习 2: \texttt{rnn\_forward}}

对 $T_x$ 个时间步循环调用 \texttt{rnn\_cell\_forward}. 实现要点:
\begin{itemize}
    \item 从 \texttt{x.shape} 获取 $n_x,m,T_x$.
    \item 根据 \texttt{parameters['Wya']} 获取 $n_y,n_a$.
    \item 初始化 \texttt{a} 和 \texttt{y\_pred} 为零.
    \item 用 \texttt{a0} 初始化 \texttt{a\_next}.
    \item 每步保存隐藏状态、预测结果和 cache.
\end{itemize}

\subsubsection*{4.1.3 LSTM 单元}

\textbf{练习 3: \texttt{lstm\_cell\_forward}}

将 $a_{\mathrm{prev}}$ 与 $x_t$ 沿行方向拼接后, 依次计算遗忘门、输入门、候选细胞状态、输出门, 最终得到 \texttt{c\_next} 和 \texttt{a\_next}. 函数返回:
\begin{align}
    \texttt{a\_next},\quad \texttt{c\_next},\quad \texttt{yt\_pred},\quad \texttt{cache}
\end{align}
注意 cache 中应保留所有门控变量, 以便选做反向传播使用.

\subsubsection*{4.1.4 完整 LSTM 前向传播}

\textbf{练习 4: \texttt{lstm\_forward}}

对 $T_x$ 个时间步循环调用 \texttt{lstm\_cell\_forward}. 需要同时维护隐藏状态 \texttt{a\_next} 与细胞状态 \texttt{c\_next}. 返回顺序为:
\begin{align}
    \texttt{a},\quad \texttt{y},\quad \texttt{c},\quad \texttt{caches}
\end{align}
这与 \texttt{rnn\_forward} 的返回值不同, 请严格按 Notebook 要求实现.

\subsubsection*{4.1.5 反向传播(选做)}

Notebook 中提供了 RNN 和 LSTM 反向传播框架. 选做内容包括:
\begin{itemize}
    \item \texttt{rnn\_cell\_backward} 与 \texttt{rnn\_backward}.
    \item \texttt{lstm\_cell\_backward} 与 \texttt{lstm\_backward}.
\end{itemize}
完成时要注意时间维度上的梯度累加, 尤其是 \texttt{da\_prev} 和 \texttt{dc\_prev} 如何从未来时间步传回当前时间步.

\subsection*{实验4.2: 恐龙名字字符级语言模型(W1A2)}

\textbf{实验目标: }训练 RNN 生成具有恐龙名字风格的新字符串.

\subsubsection*{4.2.1 数据预处理}

读取 \texttt{dinos.txt}, 统计字符表, 建立 \texttt{char\_to\_ix} 和 \texttt{ix\_to\_char}. 每个名字作为一个训练样本, 输入序列首位补 \texttt{None}, 标签序列在末尾追加换行符 \texttt{\textbackslash n}.

\subsubsection*{4.2.2 梯度裁剪}

\textbf{练习 1: \texttt{clip}}

对梯度字典中的 \texttt{dWaa}、\texttt{dWax}、\texttt{dWya}、\texttt{db}、\texttt{dby} 做原地裁剪:
\begin{align}
    g_{ij}\leftarrow \min(\max(g_{ij},-\texttt{maxValue}),\texttt{maxValue})
\end{align}
推荐使用 \texttt{np.clip(gradient, -maxValue, maxValue, out=gradient)}.

\subsubsection*{4.2.3 字符采样}

\textbf{练习 2: \texttt{sample}}

采样流程:
\begin{enumerate}
    \item 初始化输入 $x$ 和隐藏状态 $a_{\mathrm{prev}}$ 为零.
    \item 单步前向传播得到 softmax 概率.
    \item 使用 \texttt{np.random.choice} 按概率采样字符 id.
    \item 将采样字符转为 one-hot 作为下一步输入.
    \item 采样到换行符或达到最大长度时停止.
\end{enumerate}

\subsubsection*{4.2.4 优化步骤与训练}

\textbf{练习 3: \texttt{optimize}}

整合一步训练:
\begin{align}
    \text{forward}\to \text{backward}\to \text{clip}\to \text{update}
\end{align}
函数返回当前损失、裁剪后的梯度和最后隐藏状态.

\textbf{练习 4: \texttt{model}}

完成训练循环:
\begin{itemize}
    \item 循环选择一个恐龙名字.
    \item 构造 \texttt{X=[None]+name\_indices} 和 \texttt{Y=name\_indices+[newline]}.
    \item 调用 \texttt{optimize} 更新参数.
    \item 使用指数平滑损失观察训练趋势.
    \item 每隔一定迭代次数采样生成名字.
\end{itemize}

\subsection*{实验4.3: 从零实现 Karpathy minimalRNN(2025 补充实验)}

\textbf{实验目标: }把 \texttt{minimalRNN.py} 拆成可测试的课堂练习, 完整实现字符级 Vanilla RNN 的训练闭环.

\subsubsection*{4.3.1 字符词表与参数初始化}

\begin{itemize}
    \item \textbf{练习 1}: 使用 \texttt{sorted(set(data))} 构建字符表、\texttt{char\_to\_ix}、\texttt{ix\_to\_char}.
    \item \textbf{练习 2}: 初始化 \texttt{Wxh}、\texttt{Whh}、\texttt{Why}、\texttt{bh}、\texttt{by}. 权重为小随机数, 偏置为零.
\end{itemize}

\subsubsection*{4.3.2 前向传播与 BPTT}

\begin{itemize}
    \item \textbf{练习 3}: 实现 \texttt{one\_hot} 和数值稳定的 \texttt{softmax}.
    \item \textbf{练习 4}: 实现 \texttt{rnn\_forward}, 保存 \texttt{xs}、\texttt{hs}、\texttt{ys}、\texttt{ps}, 并累加负对数似然损失.
    \item \textbf{练习 5}: 实现 \texttt{rnn\_backward}. 关键是从最后一个时间步反向遍历, 用 \texttt{dhnext} 把未来梯度传回当前时间步.
\end{itemize}

\subsubsection*{4.3.3 训练稳定性与采样}

\begin{itemize}
    \item \textbf{练习 6}: 实现 \texttt{clip\_gradients} 和 \texttt{lossFun}.
    \item \textbf{练习 7}: 实现 \texttt{sample}, 从模型概率分布中逐字符采样.
    \item \textbf{练习 8}: 实现 \texttt{train\_rnn}, 包括数据指针移动、隐藏状态续接、平滑损失、AdaGrad 更新和定期采样.
\end{itemize}

完成后应比较训练前随机采样与训练后采样的差异, 并说明模型是否学到了换行、元音、常见名字片段等字符规律.


% ============================================================
\section{核心函数对照总结}
% ============================================================

\begin{center}
\begin{tabular}{|l|l|l|}
\hline
\textbf{函数} & \textbf{所在实验} & \textbf{作用} \\
\hline
\texttt{rnn\_cell\_forward} & W1A1 & 单步 RNN 前向传播 \\
\texttt{rnn\_forward} & W1A1 & 完整 RNN 前向传播 \\
\texttt{lstm\_cell\_forward} & W1A1 & 单步 LSTM 前向传播 \\
\texttt{lstm\_forward} & W1A1 & 完整 LSTM 前向传播 \\
\texttt{clip} & W1A2 & 裁剪 RNN 梯度 \\
\texttt{sample} & W1A2 / 2025 & 字符级采样生成 \\
\texttt{optimize} & W1A2 & 一步训练更新 \\
\texttt{model} & W1A2 & 恐龙名字模型训练循环 \\
\texttt{lossFun} & 2025 & minimalRNN 核心损失与梯度函数 \\
\texttt{train\_rnn} & 2025 & 有限步 minimalRNN 训练闭环 \\
\hline
\end{tabular}
\end{center}


% ============================================================
\section{实验作业}
% ============================================================
\begin{enumerate}
    \item 完成 \texttt{W1A1} 中 RNN 与 LSTM 前向传播相关函数, 并通过测试单元格.
    \item 完成 \texttt{W1A2} 中 \texttt{clip}、\texttt{sample}、\texttt{optimize}、\texttt{model}, 训练恐龙名字生成模型.
    \item 完成 \texttt{2025-实验4.ipynb}, 说明 \texttt{minimalRNN.py} 中前向传播、BPTT、梯度裁剪、AdaGrad 更新和采样生成分别对应 Notebook 中哪些函数.
    \item 记录至少 5 个模型生成的名字样本, 分析其合理性和不足.
    \item \textbf{选做}: 完成 W1A1 中 RNN/LSTM 反向传播函数.
    \item \textbf{选做}: 为 \texttt{sample} 加入 temperature 参数, 对比不同温度下生成文本的多样性.
    \item \textbf{选做}: 将 \texttt{2025-实验4.ipynb} 中的小模型迁移到更大的纯文本语料上, 分析训练损失和生成样例的变化.
\end{enumerate}


% ============================================================
\section{提交要求}
% ============================================================

请在提交前检查材料是否齐全. 本课程作业上传只接受文本文件, 不接受 PDF、Word 文档、图片、截图、压缩包等二进制文件. 采用这一规则, 一方面是为了引入 AI 辅助评阅, 另一方面也是希望大家养成用清晰文本表达实验过程和结论的习惯.

如果实验报告确实需要配图, 请优先思考是否可以用文字、公式或表格表达清楚; 若必须插图, 请将图片放入 Notebook 中, 形成一份完整的 \texttt{.ipynb} 实验报告. 若使用 \LaTeX{} 或 Word 撰写报告, 请使用 \texttt{pandoc} 等工具转换为 \texttt{.md} 或 \texttt{.txt} 等文本格式后再提交, 不要提交 \texttt{.pdf} 或 \texttt{.docx}.

\begin{enumerate}
    \item \textbf{Jupyter Notebook 文件(.ipynb)}
    \begin{enumerate}
        \item 必交内容包括本实验要求完成的 Notebook 文件, 至少包含 \texttt{W1A1}、\texttt{W1A2} 和 \texttt{2025-实验4.ipynb}.
        \item Notebook 中应包含完整代码、测试结果、训练输出、生成样例和必要分析文字.
        \item 只需提交 Notebook 源文件, 不需要额外导出为 PDF、HTML、Python 脚本、Word 文档、截图或图片文件.
    \end{enumerate}
    \item \textbf{实验过程与 AI 互动记录文本文件(.md / .txt)}
    \begin{enumerate}
        \item 必须额外提交一份文本文件, 作为实验过程与 AI 互动记录, 推荐使用 \texttt{.md} 或 \texttt{.txt}.
        \item 记录内容至少包括: 查阅了哪些资料、向 AI 提出了哪些问题、如何继续思考、调试、修改代码或完善实验分析.
        \item 不要把 AI 的回复原文放入记录中, 也不要复制粘贴大段代码; 记录应只包含你自己写的内容, 包括你的问题、判断、尝试、修改思路和必要的简短总结.
        \item 记录应尽量按时间顺序整理, 不要只保留最终答案.
    \end{enumerate}
\end{enumerate}


% ============================================================
\section{关于作业的重要说明}
% ============================================================

本实验的核心目标是理解序列模型的递推结构和训练闭环, 而不是只得到一份能运行的代码. 代码填空部分应由学生自己完成, 并通过调试和单元测试逐步确认正确性.

可以使用 AI 工具辅助理解公式、解释报错、查询 NumPy API 或检查矩阵维度, 但不得直接让 AI 代写全部必做代码后不加理解地提交. 若使用 AI 工具, 必须在文本记录中如实写明自己的提问、自己的判断和后续修改过程.

AI 互动记录应只保留学生自己写的内容. 不要粘贴 AI 的完整回复, 不要粘贴复制来的大段代码, 也不要把聊天记录原样导出后提交.

选做内容可以更自由地使用 AI 工具, 不需记录使用过程。

\end{document}
